在惯性实验室系 \(S\) 中，质量为 \(m>0\) 的理想火箭于 \(t=x=0\) 从静止开始沿 \(+x\) 方向运动，其加速度计始终读出恒定固有加速度 \(a>0\)。令火箭固有时为 \(\tau\)，真空光速为 \(c\)，忽略燃料喷射的质量变化，并将维持加速的外部装置视为理想供能装置。
(1) 从四速度的归一化及固有加速度定义出发，求快速度、实验室速度、\(t(\tau)\)、\(x(\tau)\)，核验四速度模及初始条件。
(2) 消去 \(\tau\) 求世界线双曲方程；求坐标加速度、到达实验室距离 \(D>0\) 时的实验室时间与固有时，并核对外力做功和相对论动能。
(3) 固定在原点的实验室信号源于 \(t_e\ge0\) 沿 \(+x\) 发出频率 \(\nu_0\) 的光脉冲。求追上火箭的充要条件、相遇固有时和火箭测得的频率，解释未来视界的物理意义。
(4) 火箭在固有时 \(\tau\) 沿自身瞬时静止系的 \(-x\) 方向发射频率 \(\nu_*\) 的光，求原点接收时刻 \(t_r\) 与接收频率，分别从洛伦兹频移和相邻波峰到达间隔核对，并讨论 \(\tau\to\infty\) 极限。
